\begin{center}
{\Large \bfseries \textcolor{lyred}{第十章 重积分}}
\end{center}

\vspace{6pt}
{\large \bfseries \textcolor{lyred}{第10.1 节 二重积分的概念与性质}}

\vspace{4pt}
{\bfseries \textcolor{lyred}{一.二重积分的概念}}

{\bfseries \textcolor{lyred}{1.曲顶柱体的体积}}

{\bfseries \textcolor{lyred}{2.平面薄片的质量}}

{\bfseries \textcolor{lyred}{3.二重积分的定义}}

设
(
)
,
f x y 在有界闭区域D 上有界,将D 任意分成n 小块
i
\sigma 
\Delta 
,作二重积分和 
(
)
,
n
i
i
i
i
f \psi \eta 
\sigma 
=
\Delta 
\sum 
, (
)
,
i
i
i
\psi \eta 
\sigma 
\forall 
\in \Delta 
,若在无限细分D 的过程中,随
\lambda \to 
,该和总是 
趋向于同一常数I ,它只依赖于
(
)
,
f x y 和D ,则称
(
)
,
f x y 在D 上\textcolor{lyred}{可积},记 
(
)
,
D
I
f x y d\sigma 
=\int \int 
,称它为
(
)
,
f x y 在D 上的\textcolor{lyred}{二重积分},其中
(
)
,
f x y 为\textcolor{lyred}{被积函数}, 
(
)
,
f x y d\sigma 为\textcolor{lyred}{被积表达式}, d\sigma 为\textcolor{lyred}{面积元素},在直角坐标下,也可以记d
dxdy
\sigma =
, 
D 为\textcolor{lyred}{积分区域}, ,x y 为\textcolor{lyred}{积分变量}. 
\textcolor{lyred}{定理}.平面有界闭区域上的二元连续函数必可积. 
\textcolor{lyred}{例}.
x
y
a
a
x
y d
a
\sigma 
\pi 
+
\le 
-
-
=
\int \int 
. 
\textcolor{lyred}{例}.
x
y
a
x
y d
a
a
a
a
a
\sigma 
\pi 
\pi 
\pi 
+
\le 
+
=
\cdots 
-
\cdots 
=
\int \int 
. 
\vspace{4pt}
{\bfseries \textcolor{lyred}{二.二重积分的性质}}

\textcolor{lyred}{性质1}.
(
)
(
)
(
)
(
)
,
,
,
,
D
D
D
f x y
g x y
d
f x y d
g x y d
\alpha 
\beta 
\sigma 
\alpha 
\sigma 
\beta 
\sigma 
\pm 
=
\pm 




\int \int 
\int \int 
\int \int 
; 
\textcolor{lyred}{性质2}.
(
)
(
)
(
)
,
,
,
D
D
D
D
f
x y d
f x y d
f
x y d
\sigma 
\sigma 
\sigma 
+
=
+
\int \int 
\int \int 
\int \int 
; 
\textcolor{lyred}{性质3}.
D
d\sigma 
\sigma 
\cdots 
=
\int \int 
,即区域D 的面积,一般地,
D
k d
k
\sigma 
\sigma 
\cdots 
=
\int \int 
; 
\textcolor{lyred}{性质4}.设在D 上
(
)
(
)
,
,
f x y
g x y
\le 
,则
(
)
(
)
,
,
D
D
f x y d
g x y d
\sigma 
\sigma 
\le 
\int \int 
\int \int 
; 
\textcolor{lyred}{推论}.
(
)
(
)
,
,
D
D
f x y d
f x y d
\sigma 
\sigma 
\le 
\int \int 
\int \int 
. 
\textcolor{lyred}{性质5(估值定理)}.设在D 上
(
)
,
m
f x y
M
\le 
\le 
,则
(
)
,
D
m
f x y d
M
\sigma 
\sigma 
\sigma 
\le 
\le 
\int \int 
; 
\textcolor{lyred}{性质6(积分中值定理)}.设
(
)
,
f x y 在D 上连续,则存在(
)
,
D
\psi \eta \in 
,使得 
(
)
(
)
,
,
D
f
f
x y d
\psi \eta 
\sigma 
\sigma 
=
\int \int 
,称为
(
)
,
f x y 在D 上的\textcolor{lyred}{平均值}. 
\vspace{4pt}
{\bfseries \textcolor{lyred}{三.对称性}}

1.设D 关于y 轴对称,即(
)
(
)
,
,
x y
D
x y
D
\in 
\Rightarrow -
\in 
,则 
(1)当
(
)
(
)
,
,
f
x y
f x y
-
=-
时,
(
)
,
D
f x y d\sigma =
\int \int 
; 
(2)当
(
)
(
)
,
,
f
x y
f x y
-
=
时,
(
)
(
)
\{
\}
,
,
D
D
x
f x y d
f x y d
\sigma 
\sigma 
\cap 
\ge 
=
\int \int 
\int \int 
. 
2.设D 关于x 轴对称,即(
)
(
)
,
,
x y
D
x
y
D
\in 
\Rightarrow 
-
\in 
,则情况类似. 
3.设D 关于直线y
x
=
对称,即轮换对称:(
)
(
)
,
,
x y
D
y x
D
\in 
\Rightarrow 
\in 
,则 
(
)
(
)
(
)
(
)
,
,
,
,
D
D
D
f x y d
f
y x d
f x y
f
y x
d
\sigma 
\sigma 
\sigma 
=
=
+




\int \int 
\int \int 
\int \int 
;特别地,此时 
()
()
()
()
D
D
D
f x d
f
y d
f x
f
y
d
\sigma 
\sigma 
\sigma 
=
=
+




\int \int 
\int \int 
\int \int 
. 
\textcolor{lyred}{例}.求
(
)
D
I
x
y
d\sigma 
=
+
+
\int \int 
,其中
:
x
y
D a
b
+
\le . 
解.
D
I
d
ab
ab
\sigma 
\pi 
\pi 
=
=
=
\int \int 
. 
\textcolor{lyred}{例}.求
(
)
D
I
x
y
d\sigma 
=
+
\int \int 
,其中
:
D x
y
+
\le . 
解.(
)
x
y
+
展开式里每一项均是x 或y 的奇函数,故
I =
. 
\textcolor{lyred}{例}.求
(
)
(
)
x
x
D
I
y
x
e
x
e
d\sigma 
-


=
+
+
-


\int \int 
,其中
:
D x
y
\le 
\le , 1
x
-\le 
\le . 
解.
(
)
(
)
x
x
x
x
D
D
I
xy e
e
d
y e
e
d
\sigma 
\sigma 
-
-
=
+
+
-
=
\int \int 
\int \int 
. 
\textcolor{lyred}{例}.求
D
I
xd\sigma 
=\int \int 
,其中
(
)
(
)
:
D
x
y
-
+
-
\le . 
解.
(
)
(
)
D
D
D
I
x
d
x
d
d
\sigma 
\sigma 
\sigma 
\pi 
\pi 
=
-+
=
-
+
=
+
=
\int \int 
\int \int 
\int \int 
. 
\textcolor{lyred}{例}.求
D
I
yd\sigma 
=\int \int 
,其中D 由
x
y
y
=-
-
,
x =-,
y =
,
y =
围成. 
解.
(
)
(
)
1 1
D
D
D
D
I
y
d
y
d
d
d
\pi 
\sigma 
\sigma 
\sigma 
\sigma 
=
-+
=
-
+
=
=
-
\int \int 
\int \int 
\int \int 
\int \int 
. 
\textcolor{lyred}{例}.求
(
)
D
I
x
y
dxdy
=
-
\int \int 
,其中
:
D x
y
x
y
+
\le 
+
. 
解.
D
D
I
x dxdy
y dxdy
=
-
=
\int \int 
\int \int 
,或者
(
)
D
I
y
x
dxdy
=
-
=
\int \int 
. 
\textcolor{lyred}{例}.求
()
()
()
()
D
af x
bf
y
I
d
f x
f
y
\sigma 
+
=
+
\int \int 
,其中
:
D x
y
+
\le . 
解.
()
()
()
()
()
()
()
()
D
D
af x
bf
y
af
y
bf x
a
b
a
b
I
d
d
f x
f
y
f
y
f x
\sigma 
\sigma 
\pi 


+
+
+
+
=
+
=
=


+
+


\int \int 
\int \int 
. 
\textcolor{lyred}{例}.证明:
x
y
D
e
I
d
e
\sigma 
\pi 
+
=
\ge 
+
\int \int 
,其中
:
D x
y
+
\le . 
证.若
()
f x >
,则
()
()
()
()
()
()
D
D
D
f x
f x
f
y
d
d
d
f
y
f
y
f x
\sigma 
\sigma 
\sigma 


=
+
\ge 




\int \int 
\int \int 
\int \int 
,证毕. 
\textcolor{lyred}{补充练习 }
1.设
(
)
,
f x y 连续,求
(
)
(
)
(
)
lim
,
r
x x
y y
r
f x y d
r
\sigma 
\pi 
+
\to 
-
+
-
\le 
\int \int 
. 
解.上式
(
)
(
)
(
)
(
)
(
)
,
,
lim
,
lim
,
,
x
y
r
f
f
f x
y
\psi \eta 
\psi \eta 
\psi \eta 
+
\to 
\to 
=
=
=
. 
2.设
(
)
(
)
,
3 1
,
x
y
f x y
x
y
f x y d\sigma 
+
\le 
=
-
-
+\int \int 
,求
(
)
,
f x y . 
解.设
(
)
,
D
f x y d
A
\sigma =
\int \int 
,则
(
)
,
3 1
f x y
x
y
A
=
-
-
+
,于是 
D
A
x
y d
A
A
A
\pi 
\sigma 
\pi 
\pi 
\pi 
\pi 
=
-
-
+
=
+
\Rightarrow 
=-
\int \int 
. 
3.求
()
(
)
D
I
f x
f
y
d\sigma 
=
-
-




\int \int 
,其中
:
D
x
y
+
\le . 
解.
()
(
)
()
(
)
2 D
I
f x
f
x
f
y
f
y
d\sigma 
=
-
-
+
-
-
=




\int \int 
. 
\vspace{6pt}
{\large \bfseries \textcolor{lyred}{第10.2 节 二重积分的计算法}}

\vspace{4pt}
{\bfseries \textcolor{lyred}{一.利用直角坐标计算二重积分}}

\textcolor{lyred}{公式1}.设D 为X-型区域: a
x
b
\le 
\le 
,
()
()
y
x
y
y
x
\le 
\le 
,则 
(
)
(
)
()
()
(
)
()
()
,
,
,
y
x
y
x
b
b
D
a
y
x
a
y
x
f x y d
f x y dy dx
dx
f x y dy
\sigma 


=
=






\int \int 
\int \int 
\int 
\int 
. 
\textcolor{lyred}{公式2}.设D 为Y-型区域:c
y
d
\le 
\le 
,
()
()
x
y
x
x
y
\le 
\le 
,则 
(
)
(
)
()
()
(
)
()
()
,
,
,
x
y
x
y
d
d
D
c
x
y
c
x
y
f x y d
f x y dx dy
dy
f x y dx
\sigma 


=
=






\int \int 
\int \int 
\int 
\int 
. 
\textcolor{lyred}{推论}.
()
()
[
][
]
()
()
,
,
b
d
a b
c d
a
c
x
y dxdy
x dx
y dy
\varphi 
\zeta 
\varphi 
\zeta 
\times 
\cdots 
=
\cdots 
\int \int 
\int 
\int 
. 
\textcolor{lyred}{例}.设
()
f x 连续,且
()
f x >
,证明:
()
()
(
)
b
b
a
a
f x dx
dx
b
a
f x
\ge 
-
\int 
\int 
. 
证.左式
()
()
[
][
]
()
()
()
()
[
][
]
,
,
,
,
a b
a b
a b
a b
D
f x
f x
f
y
d
d
d
f
y
f
y
f x
\sigma 
\sigma 
\sigma 
\times 
\times 


=
=
+
\ge 




\int \int 
\int \int 
\int \int 
,证毕. 
\textcolor{lyred}{例}.将
(
)
,
D
I
f x y d\sigma 
=\int \int 
化作两种顺序的二次积分,其中 
(1) D 由
y
x
=
, y
x
=
围成; 
解.
(
)
(
)
,
,
y
x
y
x
I
dx
f x y dy
dy
f x y dx
=
=
\int 
\int 
\int 
\int 
. 
(2) D 由
y
x
=
-
,
x
y
+
-
=
围成; 
解.
(
)
(
)
4 2
,
,
y
x
x
y
I
dx
f x y dy
dy
f x y dx
-
-
-
-
=
=
\int 
\int 
\int 
\int 
. 
(3) D 由
2y
x
=
,
y
x
=
-
围成; 
解.
(
)
(
)
(
)
,
,
,
y
x
x
x
x
y
I
dx
f x y dy
dx
f x y dy
dy
f x y dx
+
-
-
-
=
+
=
\int 
\int 
\int 
\int 
\int 
\int 
. 
(4) D 由
x
y
+
=
,
x
y
-
=
,
x =
围成. 
解.
(
)
(
)
(
)
,
,
,
y
y
x
x
I
dx
f x y dy
dy
f x y dx
dy
f x y dx
-
-
=
=
+
\int 
\int 
\int \int 
\int 
\int 
. 
(5) D 由
ln
y
x
=
,
y
e
x
=
+-
,
y =
围成. 
解.
(
)
(
)
(
)
ln
,
,
,
y
e
y
e
x
e
e
x
e
e
I
dx
f x y dy
dx
f x y dy
dy
f x y dx
+-
+
+-
=
+
=
\int 
\int 
\int 
\int 
\int 
\int 
. 
(6) D 由y
x
=
,
y
x
=
+,
y =,
y =
围成. 
解.
(
)
(
)
(
)
(
)
,
,
,
,
y
x
x
x
x
y
I
dx
f x y dy
dx
f x y dy
dx f x y dy
dy
f x y dx
+
+
-
=
+
+
=
\int 
\int 
\int 
\int 
\int \int 
\int 
\int 
. 
\textcolor{lyred}{例}.求
D
I
y d\sigma 
=\int \int 
,其中D 由
(
)
(
)(
)
sin
1 cos
x
a t
t
t
y
a
t
\pi 
=
-

\le \le 

=
-

与x 轴围成. 
解.
()
(
)
(
)
1 cos
sin
y x
a
a
a
I
dx
y dy
y dx
t
d t
t
a
\pi 
\pi 
\pi 
\pi 
=
=
=
-
-
=
\int 
\int 
\int 
\int 
. 
\textcolor{lyred}{例}.求
arctan
D
I
yd\sigma 
=\int \int 
,其中D 由y
x
=
,
y =,
x =
围成. 
解.
arctan
arctan
y
I
dy
ydx
y
ydy
\pi 
=
=
=
-
\int \int 
\int 
. 
\textcolor{lyred}{例}.求
y
D
I
x e
d\sigma 
-
=\int \int 
,其中D 由y
x
=
,
y =,
x =
围成. 
解.
y
y
y
y
I
dy x e
dx
e
dy
e
-
-


=
=
=
-




\int \int 
\int 
. 
\textcolor{lyred}{例}.
sin
sin
sin
sin
1 cos1
x
y
D
x
x
x
dy
dx
d
dx
dy
xdx
x
x
x
\sigma 
+
-
=
=
=
=-
\int 
\int 
\int \int 
\int 
\int 
\int 
. 
\textcolor{lyred}{例}.
x
D
y
x
dy
x
y dx
x
y d
dx
x
y dy
dx
\pi 
\pi 
\sigma 
-
=
-
=
-
=
=
\int 
\int 
\int \int 
\int 
\int 
\int 
. 
\textcolor{lyred}{例}.设
()
2x
y
x
f x
e
dy
-
=\int 
,求
()
I
f x dx
+\infty 
=\int 
. 
解.
y
x
y
y
y
y
x
I
dx
e
dy
dy e
dx
ye
dy
+\infty 
+\infty 
+\infty 
-
-
-
=
=
=
=
\int 
\int 
\int 
\int 
\int 
. 
\textcolor{lyred}{例}.交换积分顺序:(1)
(
)
(
)
,
,
y
x
dx
f x y dy
dy
f x y dx
-
-
-
=
\int 
\int 
\int 
\int 
; 
(2)
(
)
(
)
(
)
,
,
,
y
y
x
x
dx
f x y dy
dy
f x y dx
dy
f x y dx
-
-
=
+
\int 
\int 
\int 
\int 
\int 
\int 
; 
(3)
(
)
(
)
(
)
,
,
,
y
x
y
x
x
dy
f x y dx
dx
f x y dy
dx
f x y dy
=
+
\int 
\int 
\int 
\int 
\int 
\int 
; 
(4)
(
)
(
)
(
)
2 2
,
,
,
x
y
x
y
dy
f x y dx
dx
f x y dy
dx
f x y dy
-
-
=
+
\int 
\int 
\int 
\int 
\int 
\int 
; 
(5)
(
)
(
)
(
)
,
,
,
y
x x
x
y
dx
f x y dy
dx
f x y dy
dy
f x y dx
-
-
-
-
-
+
=
\int 
\int 
\int 
\int 
\int 
\int 
; 
(6)
(
)
(
)
(
)
,
,
,
y
x
x
y
y
dy
f x y dx
dy
f x y dx
dx
f x y dy
+
=
\int 
\int 
\int 
\int 
\int 
\int 
. 
\textcolor{lyred}{例}.求
D
I
y
x d\sigma 
=
-
\int \int 
,其中
:
D x \le ,0
y
\le 
\le 
. 
解.
(
)
(
)
x
x
I
dx
x
ydy
dx
y
x dy
x
dx
x
dx
-
-
-
-
=
-
+
-
=
+
-
=
\int 
\int 
\int 
\int 
\int 
\int 
(
)
2sin
2 sin
x dx
t
d
t
\pi 
\pi 
\pi 
\pi 
-
-
+
-
=
+
+
=
+
\int 
\int 
. 
\textcolor{lyred}{例}.求
(
)
cos
D
I
x
y d\sigma 
=
+
\int \int 
,其中D 由
y =
, y
x
=
,
x
\pi 
=
围成. 
解.
(
)
(
)
cos
cos
y
x
y
x
I
dy
x
y dx
dx
x
y dy
\pi 
\pi 
\pi 
\pi 
\pi 
\pi 
-
-
=
+
-
+
=
-
\int 
\int 
\int 
\int 
. 
\textcolor{lyred}{例}.求
(
)
sin
D
I
xy
y d\sigma 
=
+
\int \int 
,其中D 由x
\pi 
=-
, y
\pi 
=
, y
x
=
围成. 
解.
(
)
sin
sin
1 cos
D
x
I
ydxdy
dx
ydy
x dx
\pi 
\pi 
\pi 
\pi 
=
=
=
+
=
\int \int 
\int \int 
\int 
. 
\textcolor{lyred}{例}.求
(
)
D
I
y
xy
x
y
d\sigma 
=
+
\cdots 
+
\int \int 
,其中D 由
y
x
=
,
y =,
x =-围成. 
解.
D
D
I
y d
dx y dy
\sigma 
+
-
=
=
=
\int \int 
\int 
\int 
. 
\vspace{4pt}
{\bfseries \textcolor{lyred}{二.利用极坐标计算二重积分}}

\textcolor{lyred}{公式3}.设D 为曲边扇形:\alpha 
\theta 
\beta 
\le 
\le 
,
()
()
\rho \theta 
\rho 
\rho 
\theta 
\le 
\le 
,则 
(
)
(
)
()
()
,
cos ,
sin
D
f x y d
d
f
d
\rho 
\theta 
\beta 
\alpha 
\rho \theta 
\sigma 
\theta 
\rho 
\theta \rho 
\theta \rho \rho 
=
\int \int 
\int 
\int 
. 
\textcolor{lyred}{例}.求
z
x
y
=
+
与
z
x
y
=
-
-
所围立体的体积. 
解.
(
)(
)
(
)
x
y
x
y
V
x
y
x
y
d
x
y
d
\sigma 
\sigma 
+
\le 
+
\le 


=
-
-
-
+
=
-
-
=


\int \int 
\int \int 
(
)
(
)
d
d
d
\pi 
\theta 
\rho 
\rho \rho 
\pi 
\rho 
\rho \rho 
\pi 
-
=
-
=
\int 
\int 
\int 
. 
\textcolor{lyred}{例}.求
(
)
D
I
x
y
d\sigma 
=
+
\int \int 
,其中
:1
D
x
y
\le 
+
\le 
. 
解.
(
)
(
)
D
D
I
x
y
d
x
y
d
d
d
\pi 
\pi 
\sigma 
\sigma 
\theta \rho 
\rho \rho 
=
+
=
+
=
\cdots 
=
\int \int 
\int \int 
\int 
\int 
. 
\textcolor{lyred}{例}.
(
)
x
y
R
x
y
R
x
y
R
x
y
x
y
d
d
x
y
d
a
b
a
b
a
b
\sigma 
\sigma 
\sigma 
+
\le 
+
\le 
+
\le 






+
=
+
=
+
+












\int \int 
\int \int 
\int \int 
R
a
b
\pi 


=
+




. 
\textcolor{lyred}{例}.求
(
)
D
I
xy
x x
y
y
x
y
d\sigma 
=
+
+
+
+
\int \int 
,其中
:
D x
y
ax
+
\le 
. 
解.
2 cos
cos
cos
a
D
I
x x
y d
d
d
a
d
a
\pi 
\pi 
\theta 
\pi 
\pi 
\sigma 
\theta 
\rho 
\theta \rho \rho \rho 
\theta \theta 
-
-
=
+
=
\cdots 
\cdots 
=
=
\int \int 
\int 
\int 
\int 
. 
\textcolor{lyred}{例}.求
x
y
z
+
+
\le 
被
x
y
x
+
=
截得立体的体积. 
解.
2cos
x
y
x
V
x
y dxdy
d
d
\pi 
\theta 
\pi 
\pi 
\theta 
\rho 
\rho \rho 
+
\le 
-


=
-
-
=
-
\cdots 
=
-




\int \int 
\int 
\int 
. 
\textcolor{lyred}{例}.求
(
)
D
I
x
x yd\sigma 
=
+
\int \int 
,其中
:
D ay
x
y
ay
\le 
+
\le 
. 
解.
2 sin
sin
cos sin
sin
cos
a
D
a
I
xyd
d
d
a
d
a
\pi 
\pi 
\theta 
\theta 
\sigma 
\theta 
\rho 
\theta 
\theta \rho \rho 
\theta 
\theta \theta 
=
=
\cdots 
=
=
\int \int 
\int 
\int 
\int 
. 
\textcolor{lyred}{例}.求
D
I
x
y d\sigma 
=
-
-
\int \int 
,其中
:
D x
y
+
\le ,
x
y
y
+
\le 
. 
解.
2sin
2 3
I
d
d
d
d
\pi 
\pi 
\theta 
\pi 
\theta 
\rho \rho \rho 
\theta 
\rho \rho \rho 
\pi 
-
=
-
+
-
=
-
\int 
\int 
\int 
\int 
. 
\textcolor{lyred}{例}.求
D
I
xyd\sigma 
=\int \int 
,其中D 由
y =
,
y =
,
x =-,
x
y
y
=-
-
围成. 
解.
2sin
cos sin
I
dx xydy
d
d
\pi 
\theta 
\pi 
\theta 
\rho 
\theta 
\theta \rho \rho 
-


=
-
\cdots 
=---
=-




\int 
\int 
\int 
\int 
. 
\textcolor{lyred}{例}.求
D
I
x
y d\sigma 
=
+
\int \int 
,其中
: 2
D
x
x
y
-
\le 
+
\le 
. 
解.
2cos
I
d
d
d
d
\pi 
\pi 
\theta 
\pi 
\theta \rho \rho \rho 
\theta 
\rho \rho \rho 
\pi 
-
=
\cdots 
-
\cdots 
=
-
\int 
\int 
\int 
\int 
. 
\textcolor{lyred}{例}.求
D
I
x
y
y d\sigma 
=
+
+
\int \int 
,其中
:
D x
y
+
\le 
. 
解.
(
)
(
)
y x
y
x
y
y
I
x
y
y d
x
y
y d
\sigma 
\sigma 
-
\le 
+
\le 
+
\le -
=
+
+
-
+
+
=
\int \int 
\int \int 
(
)
(
)
x
y
x
y
y
x
y
y d
x
y
y d
\sigma 
\sigma 
\pi 
+
\le 
+
\le -
+
+
-
+
+
=
\int \int 
\int \int 
. 
\textcolor{lyred}{例}.求
arctan
D
y
I
d
x \sigma 
=\int \int 
,其中D 由y
x
=
,
y =
,
x =
围成. 
解.
cos
tan
ln 2
cos
I
d
d
d
d
\pi 
\pi 
\pi 
\theta 
\pi 
\theta 
\theta \rho \rho 
\theta 
\theta 
\theta 
\theta 
\theta 
=
\cdots 
=
\cdots 
=
=
-
\int 
\int 
\int 
\int 
. 
\textcolor{lyred}{例}.求
D
I
x
y d\sigma 
=
+
\int \int 
,其中D 由y
x
=
,
y
x
=
围成. 
解.
(
)
sin
cos
sin
tan
sec
sec
cos
I
d
d
d
d
\theta 
\pi 
\pi 
\pi 
\theta 
\theta 
\theta 
\rho \rho \rho 
\theta 
\theta 
\theta 
\theta 
\theta 
=
\cdots 
=
=
=
+
\int 
\int 
\int 
\int 
. 
\textcolor{lyred}{例}.求
y x
y x
D
I
e
d\sigma 
-
+
=\int \int 
,其中D 由
x
y
+
=与坐标轴围成. 
解.
(
)
sin
cos
sin
cos
cos
sin
sin
cos
sin
cos
cos
sin
I
d
e
d
e
d
\pi 
\pi 
\theta 
\theta 
\theta 
\theta 
\theta 
\theta 
\theta 
\theta 
\theta 
\theta 
\theta 
\rho \rho 
\theta 
\theta 
\theta 
-
-
+
+
+
=
=
=
+
\int 
\int 
\int 
(
)
(
)
tan
tan
tan
tan
u
u
u
d
du
e
e
e
d
e
u
e
u
\pi 
\theta 
\theta 
\theta 
\theta 
+\infty 
+\infty 
-
-
-
+
+
+


=
=-
=
-


+


+
+
\int 
\int 
\int 
. 
\textcolor{lyred}{例}.计算概率积分
x
I
e
dx
\pi 
+\infty 
-
-\infty 
=
\int 
. 
解.
x
x
x
y
I
e
dx
e
dx
dx
e
dy
d
e
d
\pi 
\rho 
\theta 
\rho \rho 
\pi 
\pi 
\pi 
+\infty 
+\infty 
+\infty 
+\infty 
+\infty 
+
-
-
-
-
-\infty 
-\infty 
-\infty 
-\infty 
=
\cdots 
=
=
=
\int 
\int 
\int 
\int 
\int 
\int 
t
e
d
e dt
\rho 
\rho 
+\infty 
+\infty 
-
-
=
=
\int 
\int 
,故
I =. 
\vspace{4pt}
{\bfseries \textcolor{lyred}{三.化二次积分为定积分}}

\textcolor{lyred}{例}.设
()
sin
f x
x
=
,令
()
()
t
t
y
F t
dy
f x dx
=\int \int 
,求
()
F'
. 
解.
()
()
()(
)
t
x
t
F t
dx
f x dy
f x
x
dx
=
=
-
\int \int 
\int 
,故
()
()
sin16
F
f
'
=
=
. 
\textcolor{lyred}{例}.设
()
f x 连续,证明:
(
)
()
(
)
()
y
b
b
n
n
a
a
a
dy
y
x
f x dx
b
x
f x dx
n
+
-
=
-
+
\int \int 
\int 
. 
证.左式
(
)
()
(
)
()
b
b
b
n
n
a
x
a
dx
y
x
f x dy
b
x
f x dx
n
+
=
-
=
-
+
\int \int 
\int 
,证毕. 
\textcolor{lyred}{例}.设
()
f x 连续,证明:
()
()
()
x
dx
f x f
y dy
f x dx


=




\int \int 
\int 
. 
证.左式
()
()
()
()
[
][
]
0,1
0,1
y
dy
f
y f x dx
f x f
y dxdy
\times 
=
=
\int \int 
\int \int 
,即得,证毕. 
\textcolor{lyred}{例}.设
()
f u 连续,证明:
(
)
()
x
y
f x
y dxdy
f u du
+
\le 
-
+
=
\int \int 
\int 
. 
证. D 由
x
y
+
=\pm ,
x
y
-
=\pm 围成,令u
x
y
v
x
y
=
+

=
-

,
(
)
(
)
,
,
x y
u v
\partial 
=
\partial 
,则 
(
)
()
()
D
f x
y dxdy
du
f u dv
f u du
-
-
-
+
=
=
\int \int 
\int 
\int 
\int 
,证毕. 
\textcolor{lyred}{补充练习 }
1.求
(
)
x
y
D
I
y
xe
d\sigma 
+
=
+
\int \int 
,其中D 由y
x
=
,
y =-,
x =围成. 
解.
y
D
I
yd
dy
ydx
y dy
\sigma 
-
-
-
=
=
=-
=-
\int \int 
\int 
\int 
\int 
. 
2.设
(
)
,
,
x
y
e
x
y
f x y
+

\le 
+
\le 

=

，
其它
,求
(
)
2,0
,
x
y
I
f x y d\sigma 
\le \le 
\le \le 
=
\int \int 
. 
解.
(
)
x
y
x
y
x
y
x
y
x
y
x
y
I
d
d
e
d
e
d
\pi 
\sigma 
\sigma 
\sigma 
\pi 
\sigma 
+
+
+
\le 
+
\ge 
\le 
+
\le 
\le 
+
\le 
=
+
+
=
+
-
+
=
\int \int 
\int \int 
\int \int 
\int \int 
(
)
(
)
(
)
e
e
d
e
d
e
e
\pi 
\rho 
\pi 
\pi 
\pi 
\pi 
\theta 
\rho \rho 
\pi 
\pi 
--
+
-
+
=
+
-
+
-
=
+
\int 
\int 
. 
3.求
D
I
x
y
d\sigma 
=
+
-
\int \int 
,其中
:
D
x \le ,
y \le . 
解.
(
)
(
)
D
x
y
x
y
I
x
y
d
x
y
d
x
y
d
\sigma 
\sigma 
\sigma 
+
\le 
+
\ge 
=
+
-
=
-
-
+
+
-
=
\int \int 
\int \int 
\int \int 
(
)
(
)
x
d
d
dx
x
y
dy
\pi 
\pi 
\pi 
\theta 
\rho 
\rho \rho 
\pi 
-


-
+
+
-
=
+
-
=
-




\int 
\int 
\int 
\int 
; 
或者,
(
)
(
)
cos
I
d
d
d
d
\pi 
\pi 
\theta 
\theta 
\rho 
\rho \rho 
\theta 
\rho 
\rho \rho 
\pi 
=
-
+
-
=
-
\int 
\int 
\int 
\int 
. 
4.求
(
)
D
I
d
x
y
\sigma 
=
+
+
\int \int 
,其中
: 0
D
x
\le 
\le ,0
y
\le 
\le . 
解.
(
)
(
)
cos
sin
I
d
d
d
d
\pi 
\pi 
\theta 
\theta 
\pi 
\theta 
\rho \rho 
\theta 
\rho \rho 
\rho 
\rho 
=
+
+
+
\int 
\int 
\int 
\int 
arctan
arctan
2 2
2 2
\pi 


=
+
-




; 
或者,
(
)
cos
arctan
I
d
d
\pi 
\theta 
\theta 
\rho \rho 
\rho 
=
=
+
\int 
\int 
. 
5.设
(
)
(
)
(
)
,
,
,
D
f x y
x
f x y d
y
f x x dx
\sigma 
-
=
+
+
\int \int 
\int 
,其中D 由y
x
=
与
y =围成,求 
(
)
,
f x y . 
解.设
(
)
,
D
A
f x y d\sigma 
=\int \int 
,
(
)
,
B
f x x dx
-
=\int 
,则
(
)
,
f x y
Ax
By
=
+
+,于是 
(
)
D
D
A
Ax
By
d
B
yd
B
\sigma 
\sigma 
=
+
+
=
+=
+
\int \int 
\int \int 
,
(
)
B
Ax
Bx
dx
-
=
+
+
=
\int 
,故 
故
A =
,因此
(
)
,
f x y
x
y
=
+
+. 
6.设
()
t
x
y
t
f t
e
f
x
y
d
\pi 
\sigma 
+
\le 


=
+
+




\int \int 
,求
()
f t . 
解.
()
t
t
t
t
f t
e
d
f
d
e
f
d
\pi 
\pi 
\pi 
\theta 
\rho \rho \rho 
\pi 
\rho \rho \rho 




=
+
=
+








\int 
\int 
\int 
,两边求导, 
()
()
()
()
t
t
f
t
te
tf t
f
t
tf t
te
\pi 
\pi 
\pi 
\pi 
\pi 
\pi 
'
'
=
+
\Rightarrow 
-
=
,又
()
f
=,解得 
()
(
)
t
f t
e
t
\pi 
\pi 
=
+
. 
7.设
(
)
,
f x y 连续,且
(
)
0,0
f
=,求
(
)
lim
,
x
x t
x
I
dt
f t u du
x
+
-
\to 
=
\int 
\int 
. 
解.
(
)
(
)
(
)
,
,
0,0
lim
lim
D
x
x
f t u dtdu
f
x
f
I
x
x
\psi \eta 
+
+
\to 
\to 
\cdots 
=
=
=
=
\int \int 
. 
8.设
(
)
,
f x y 连续,且
(
)
0,0
f
=,求
(
)
lim
,
x
tx t
x
I
dt
f t u du
x
+
-
\to 
=
\int 
\int 
. 
解.
(
)
(
)
2 2
,
,
lim
lim
,
t
u
tx u
x
x
f t u dtdu
I
f
x
\pi 
\pi 
\psi \eta 
+
+
+
\le 
\ge 
\to 
\to 
=
=
=
\int \int 
. 
9.设
(
)
,
f x y 连续,且
(
)
0,0
f
=,求
(
)
lim
,
x
x
x
t
I
dt
f t u du
x
+
\to 
=
\int \int 
. 
解.
(
)
(
)
(
)
(
)
,
,
,
0,0
lim
lim
lim
x
u
x
x
x
x
du
f t u dt
f t x dt
f
x
f
I
x
x
\psi 
+
+
+
\to 
\to 
\to 
=
=
=
=
=
\int 
\int 
\int 
. 
10.证明:
x
e
dx
e
\pi 
-




>
-








\int 
. 
证.
[
][
]
0,1
0,1
1,
0,
x
x
y
x
y
x
y
x
y
x
y
e
dx
e
dx e
dy
e
d
e
d
\sigma 
\sigma 
-
-
-
-
-
-
-
\times 
+
\le 
\ge 
\ge 

=
=
>
=




\int 
\int 
\int 
\int \int 
\int \int 
d
e
d
e
\pi 
\rho 
\pi 
\theta 
\rho \rho 
-


=
-




\int 
\int 
,证毕. 
11.设
()
f x 连续,证明:
()
()
b
b
a
a
f
x dx
f x dx
b
a


\ge 


-


\int 
\int 
. 
证.(
)
()
()
()
[
][
]
()
[
][
]
,
,
,
,
b
b
b
a
a
a
a b
a b
a b
a b
b
a
f
x dx
dy
f
x dx
f
x d
f
y d
\sigma 
\sigma 
\times 
\times 
-
=
=
=
=
\int 
\int \int 
\int \int 
\int \int 
()
()
[
][
]
()
()
[
][
]
()
,
,
,
,
b
a b
a b
a b
a b
a
f
x
f
y
d
f x f
y d
f x dx
\sigma 
\sigma 
\times 
\times 




+
\ge 
=





\int \int 
\int \int 
\int 
,证毕. 
12.设()
p x ,
()
f x , ()
[
]
,
g x
C a b
\in 
,且()
p x \ge 
,
()
f x , ()
g x 单增,证明: 
()
()
()()
()
()()()
b
b
b
b
a
a
a
a
p x f x dx p x g x dx
p x dx
f x p x g x dx
\le 
\int 
\int 
\int 
\int 
. 
证.即证
()
()()()
[
][
]
()
()()()
[
][
]
,
,
,
,
a b
a b
a b
a b
p x f x p y g y d
p x f
y p y g y d
\sigma 
\sigma 
\times 
\times 
\le 
\int \int 
\int \int 
, 
左式-右式
()()()
()
()
[
][
]
,
,
a b
a b
p x p y g y
f x
f
y
d\sigma 
\times 
=
-
=




\int \int 
()()()
()
()
[
][
]
,
,
a b
a b
p y p x g x
f
y
f x
d\sigma 
\times 
-
=




\int \int 
()()
()
()
()
()
[
][
]
,
,
2 a b
a b
p x p y
g y
g x
f x
f
y
d\sigma 
\times 
-
-
\le 






\int \int 
,证毕. 
13.设
()
f x >
,连续,单调减少,证明:
()
()
()
()
xf
x dx
f
x dx
xf x dx
f x dx
\le 
\int 
\int 
\int 
\int 
. 
证.即证
()
()
[
][
]
()
()
[
][
]
0,1
0,1
0,1
0,1
xf
x f
y d
yf
x f
y d
\sigma 
\sigma 
\times 
\times 
\le 
\int \int 
\int \int 
, 
左式-右式
(
)
()
()
[
][
]
(
)
()
()
[
][
]
0,1
0,1
0,1
0,1
x
y f
x f
y d
y
x f
y f x d
\sigma 
\sigma 
\times 
\times 
=
-
=
-
=
\int \int 
\int \int 
(
)
()
()
()
()
[
][
]
0,1
0,1
x
y
f x
f
y
f
y f x d\sigma 
\times 
-
-
\le 




\int \int 
,证毕. 
\vspace{6pt}
{\large \bfseries \textcolor{lyred}{第10.3 节 三重积分}}

\vspace{4pt}
{\bfseries \textcolor{lyred}{一.三重积分的定义}}

设
(
)
, ,
f x y z 为有界闭区域\Omega 上的有界函数,将\Omega 任意分成n 小块
iv
\Delta ,直径为
i\lambda , 
作和
(
)
,
,
n
i
i
i
i
i
f
v
\psi \eta 
=
\Delta 
\sum 
,
(
)
,
,
i
i
i
iv
\psi \eta 
\forall 
\in \Delta ,其中
iv
\Delta 表示体积,若在无限细分\Omega 的 
过程中,随着
\{\}
1max
i
i n
\lambda 
\lambda 
\le \le 
=
\to 
,该积分和总是趋向于同一个常数I ,它只依赖于 
(
)
, ,
f x y z 和\Omega ,则称
(
)
, ,
f x y z 在\Omega 上\textcolor{lyred}{可积},记
(
)
, ,
I
f x y z dv
\Omega 
=\int \int \int 
,称为
(
)
, ,
f x y z 
在\Omega 上的\textcolor{lyred}{三重积分}. 
\textcolor{lyred}{注}. dv 为\textcolor{lyred}{体积元素},在直角坐标系下也记dv
dxdydz
=
. 
\textcolor{lyred}{定理}.设
(
)
, ,
f x y z 在有界闭区域\Omega 上连续,则
(
)
, ,
f x y z dv
\Omega \int \int \int 
存在. 
\textcolor{lyred}{物理意义}.设\Omega 具有连续密度
(
)
, ,
f x y z ,则
(
)
, ,
M
f x y z dv
\Omega 
=\int \int \int 
. 
\vspace{4pt}
{\bfseries \textcolor{lyred}{二.三重积分的性质}}

\textcolor{lyred}{性质1}.
(
)
f
g dv
f dv
g dv
\alpha 
\beta 
\alpha 
\beta 
\Omega 
\Omega 
\Omega 
\pm 
=
\cdots 
\pm 
\cdots 
\int \int \int 
\int \int \int 
\int \int \int 
; 
\textcolor{lyred}{性质2}.
(
)
(
)
(
)
, ,
, ,
, ,
f x y z dv
f x y z dv
f x y z dv
\Omega +\Omega 
\Omega 
\Omega 
=
+
\int \int \int 
\int \int \int 
\int \int \int 
; 
\textcolor{lyred}{性质3}.
1 dv
\Omega 
\cdots 
=\Omega 
\int \int \int 
,一般地,
k dv
k
\Omega 
\cdots 
=\Omega 
\int \int \int 
; 
\textcolor{lyred}{性质4}.设\Omega 上
(
)
(
)
, ,
, ,
f x y z
g x y z
\le 
,则
(
)
(
)
, ,
, ,
f x y z dv
g x y z dv
\Omega 
\Omega 
\le 
\int \int \int 
\int \int \int 
; 
\textcolor{lyred}{推论}.
(
)
(
)
, ,
, ,
f x y z dv
f x y z dv
\Omega 
\Omega 
\le 
\int \int \int 
\int \int \int 
. 
\textcolor{lyred}{性质5(估值定理)}.设\Omega 上m
f
M
\le 
\le 
,则m
f dv
M
\Omega 
\cdots \Omega \le 
\le 
\cdots \Omega 
\int \int \int 
; 
\textcolor{lyred}{性质6(积分中值定理)}.设
(
)
, ,
f x y z 在\Omega 上连续,则(
)
, ,
\psi \eta 
\exists 
\in \Omega ,使得 
(
)
(
)
, ,
, ,
f
f x y z dv
\psi \eta 
\Omega 
=\Omega \int \int \int 
,称为
(
)
, ,
f x y z 在\Omega 上的\textcolor{lyred}{平均值}. 
\vspace{4pt}
{\bfseries \textcolor{lyred}{三.对称性}}

1.若\Omega 关于xOy 面(上下)对称,则当
(
)
, ,
f x y z 关于z 为奇函数时, 
(
)
, ,
f x y z dv
\Omega 
=
\int \int \int 
;关于z 为偶函数时,
(
)
(
)
\{
\}
, ,
, ,
z
f x y z dv
f x y z dv
\Omega 
\Omega \cap 
\ge 
=
\int \int \int 
\int \int \int 
. 
2.若\Omega 关于yOz 面(前后)对称,则当
(
)
, ,
f x y z 关于x 为奇函数时, 
(
)
, ,
f x y z dv
\Omega 
=
\int \int \int 
;关于x 为偶函数时,
(
)
(
)
\{
\}
, ,
, ,
x
f x y z dv
f x y z dv
\Omega 
\Omega \cap 
\ge 
=
\int \int \int 
\int \int \int 
. 
3.若\Omega 关于xOz 面(左右)对称,则当
(
)
, ,
f x y z 关于y 为奇函数时, 
(
)
, ,
f x y z dv
\Omega 
=
\int \int \int 
;关于y 为偶函数时,
(
)
(
)
\{
\}
, ,
, ,
y
f x y z dv
f x y z dv
\Omega 
\Omega \cap 
\ge 
=
\int \int \int 
\int \int \int 
. 
4.若\Omega 关于x , y 具有轮换对称性,即当(
)
, ,
x y z \in \Omega 时,(
)
, ,
y x z \in \Omega ,则 
(
)
(
)
, ,
, ,
f x y z dv
f
y x z dv
\Omega 
\Omega 
=
\int \int \int 
\int \int \int 
,特别地,
()
()
f x dv
f
y dv
\Omega 
\Omega 
=
\int \int \int 
\int \int \int 
. 
5.若\Omega 关于y , z 具有轮换对称性,则
()
()
f
y dv
f z dv
\Omega 
\Omega 
=
\int \int \int 
\int \int \int 
. 
6.若\Omega 关于x , z 具有轮换对称性,则
()
()
f x dv
f z dv
\Omega 
\Omega 
=
\int \int \int 
\int \int \int 
. 
7.若\Omega 关于x , y , z 具有轮换对称性(不变性),即当(
)
, ,
x y z \in \Omega 时,其所有轮换 
(
)
,
,
x y z
'
'
'\in \Omega ,则
()
()
()
f x dv
f
y dv
f z dv
\Omega 
\Omega 
\Omega 
=
=
\int \int \int 
\int \int \int 
\int \int \int 
. 
\textcolor{lyred}{例}.设
: x
y
z
a
\Omega 
+
+
\le 
,记
\Omega 为\Omega 位于第一卦限的部分,则有 
(
)
(
)
(
)
x
y
z
dv
x
y
z
dv
x
y
z
dv
x dv
\Omega 
\Omega 
\Omega 
\Omega 
+
+
=
+
+
=
+
+
=
\int \int \int 
\int \int \int 
\int \int \int 
\int \int \int 
. 
\textcolor{lyred}{例}.求
(
)
I
x
y
z dv
\Omega 
=
+
+
\int \int \int 
,其中
(
)
(
)
:
x
y
z
\Omega 
-
+
-
+
\le . 
解.
(
)
(
)
1 1
I
xdv
ydv
x
dv
y
dv
dv
\pi 
\Omega 
\Omega 
\Omega 
\Omega 
\Omega 
=
+
=
-
+
+
-+
=
=
\int \int \int 
\int \int \int 
\int \int \int 
\int \int \int 
\int \int \int 
. 
\textcolor{lyred}{例}.求
(
)
I
ax
by
cz dv
\Omega 
=
+
+
\int \int \int 
,其中
:
z
x
y
z
z
\Omega 
\le 
+
+
\le 
. 
解.
I
c
zdv
c
zdv
c
zdv
c
dv
c
dv
c
\pi 
\Omega 
\Omega 
\Omega 
\Omega 
\Omega 
=
=
-
=
-
=
\int \int \int 
\int \int \int 
\int \int \int 
\int \int \int 
\int \int \int 
. 
\vspace{4pt}
{\bfseries \textcolor{lyred}{四.直角坐标下的计算}}

{\bfseries \textcolor{lyred}{1.投影法(坐标面投影法,先一后二)}}

设
(
)
(
)
(
)(
)
\{
\}
, ,
:
,
,
,
,
xy
x y z
z
x y
z
z
x y
x y
D
\Omega =
\le 
\le 
\in 
为XY-型区域,则 
(
)
(
)
(
)
(
)
(
)
(
)
(
)
,
,
,
,
, ,
, ,
, ,
xy
xy
z
x y
z
x y
D
z
x y
D
z
x y
f x y z dv
f x y z dz dxdy
dxdy
f x y z dz
\Omega 


=
=






\int \int \int 
\int \int 
\int 
\int \int 
\int 
. 
\textcolor{lyred}{例}.求I
zdv
\Omega 
=\int \int \int 
,其中\Omega 由
x =,
y =,
z
x
y
=
+
和坐标面围成. 
解.
(
)
1,0
x
y
x
y
x
y
I
dxdy
zdz
dx dy
zdz
dx
x
y
dy
+
+
\le \le 
\le \le 
=
=
=
+
=
\int \int 
\int 
\int \int 
\int 
\int \int 
. 
\textcolor{lyred}{例}.求I
zdv
\Omega 
=\int \int \int 
,其中\Omega 由
x
y
z
+
+
=和三个坐标面围成. 
解.
(
)
xy
x
x
x
y
x
y
D
I
dxdy
xdz
dx
dy
xdz
dx
x
x
y dy
-
-
--
--
=
=
=
-
-
=
\int \int 
\int 
\int 
\int 
\int 
\int 
\int 
. 
\textcolor{lyred}{例}.求
(
)
I
x
y
dv
\Omega 
=
+
\int \int \int 
,其中
:
x
y
z
\Omega 
+
\le 
, 2
z
\le 
\le . 
解.
(
)
(
)
x
y
x
y
x
y
x
y
I
dxdy
x
y
dz
dxdy
x
y
dz
+
\le 
+
+
\le 
+
=
+
-
+
=
\int \int 
\int 
\int \int 
\int 
d
d
dz
d
d
dz
\pi 
\pi 
\rho 
\rho 
\theta \rho \rho 
\rho 
\theta \rho \rho 
\rho 
\pi 
-
=
\int 
\int 
\int 
\int 
\int 
\int 
; 
或者,
(
)
(
)
x
y
x
y
x
y
I
dxdy
x
y
dz
dxdy
x
y
dz
+
\le 
\le 
+
\le 
+
=
+
+
+
=
\int \int 
\int 
\int \int 
\int 
d
d
dz
d
d
dz
\pi 
\pi 
\rho 
\theta \rho \rho \rho 
\theta \rho \rho 
\rho 
\pi 
+
=
\int 
\int 
\int 
\int 
\int 
\int 
. 
\textcolor{lyred}{例}.求
x dv
I
x
y
\Omega 
=
+
\int \int \int 
,其中\Omega 由
z =,
z
x
y
=
+
,
z
x
y
=
+
围成. 
解.
x
y
x
y
x
y
x
y
I
dxdy
x
y dz
dxdy
x
y dz
+
\le 
+
\le 
+
+
=
+
-
+
=
\int \int 
\int 
\int \int 
\int 
d
d
dz
d
d
dz
\pi 
\pi 
\rho 
\rho 
\theta \rho \rho \rho 
\theta \rho \rho \rho 
\pi 
-
=
\int 
\int 
\int 
\int 
\int 
\int 
; 
或者,
x
y
x
y
x
y
x
y
x
y
I
dxdy
x
y dz
dxdy
x
y dz
+
+
\le 
\le 
+
\le 
+
+
=
+
+
+
=
\int \int 
\int 
\int \int 
\int 
d
d
dz
d
d
dz
\rho 
\pi 
\pi 
\rho 
\rho 
\theta \rho \rho \rho 
\theta \rho \rho \rho 
\pi 
+
=
\int 
\int 
\int 
\int 
\int 
\int 
. 
\textcolor{lyred}{例}.求
(
)
I
x
y
dv
\Omega 
=
+
\int \int \int 
,其中\Omega 由
z
x
y
=
+
,
z
x
=
-
围成. 
解.
(
)
(
)(
)
x
x
y
x
y
x
y
I
dxdy
x
y
dz
x
y
x
y
dxdy
\pi 
-
+
\le 
+
+
\le 
=
+
=
+
-
-
=
\int \int 
\int 
\int \int 
. 
\textcolor{lyred}{例}.求
sin
y
z
I
dv
z
\Omega 
=\int \int \int 
,其中\Omega 由y
z
=
,
x
z
\pi 
+
=
,
x =
,
y =
围成. 
解.
sin
sin
yz
z
z
z
D
y
z
y
z
I
dydz
dx
dz
dy
dx
z
z
\pi 
\pi 
\pi 
\pi 
-
-
=
=
=
-
\int \int 
\int 
\int 
\int 
\int 
; 
或者,
sin
sin
xz
x
z
z
D
y
z
y
z
I
dxdz
dy
dx
dz
dy
z
z
\pi 
\pi 
\pi 
-
=
=
=
-
\int \int 
\int 
\int 
\int 
\int 
. 
{\bfseries \textcolor{lyred}{2.截面法(坐标轴投影法,先二后一)}}

设
(
)
(
)
\{
\}
, ,
:
,
,
z
x y z
z
z
z
x y
D
\Omega =
\le 
\le 
\in 
,则 
(
)
(
)
(
)
, ,
, ,
, ,
z
z
z
z
z
D
z
D
f x y z dv
f x y z dxdy dz
dz
f x y z dxdy
\Omega 


=
=






\int \int \int 
\int \int \int 
\int \int \int 
. 
\textcolor{lyred}{例}.求
ln
z
x
y
=
+
,
z =
,
z =围成立体的体积. 
解.
z
z
x
y
e
e
V
dv
dz
dxdy
e dz
\pi 
\pi 
\Omega 
+
\le 
-
=
\cdots 
=
\cdots 
=
=
\int \int \int 
\int 
\int \int 
\int 
. 
\textcolor{lyred}{例}.求
(
)
I
x
y
dv
\Omega 
=
+
\int \int \int 
,其中
:
x
y
z
\Omega 
+
\le 
, 2
z
\le 
\le . 
解.
(
)
z
x
y
z
I
dz
x
y
dxdy
dz
d
d
\pi 
\theta 
\rho 
\rho \rho 
\pi 
+
\le 
=
+
=
\cdots 
=
\int 
\int \int 
\int 
\int 
\int 
. 
\textcolor{lyred}{例}.求
I
x dv
\Omega 
=\int \int \int 
,其中\Omega 由
z =,
z
x
y
=
+
,
z
x
y
=
+
围成. 
解.
(
)
z
z
z
x
y
z
I
dz
x
y
d
dz
d
d
\pi 
\sigma 
\theta 
\rho 
\rho \rho 
\pi 
\le 
+
\le 
=
+
=
\cdots 
=
\int 
\int \int 
\int 
\int 
\int 
. 
\textcolor{lyred}{例}.求
(
)
I
x
y
z dv
\Omega 
=
+
+
\int \int \int 
,其中\Omega 由
x
y
z
+
+
=
和三个坐标面围成. 
解.
(
)
,
0,
x y
z x
y
I
zdv
dz
zdxdy
z
z
dz
\Omega 
+\le -
\ge 
\ge 
=
=
=
\cdots 
-
=
\int \int \int 
\int 
\int \int 
\int 
. 
\textcolor{lyred}{例}.求
I
z dv
\Omega 
=\int \int \int 
,其中
:
x
y
z
a
b
c
\Omega 
+
+
\le . 
解.
c
c
c
c
x
y
z
a
b
c
z
I
dz
z dxdy
ab
z dz
abc
c
\pi 
\pi 
-
-
+
\le -


=
=
-
=




\int 
\int \int 
\int 
. 
\textcolor{lyred}{例}.求
I
z
x
y
dv
\Omega 
=
-
+
\int \int \int 
,其中\Omega 由
x
y
+
=,
z =
,
z =围成. 
解.
(
)
(
)
x
y
z
z
x
y
I
dz
z
x
y
dxdy
dz
x
y
z dxdy
+
\le 
\le 
+
\le 
=
-
+
+
+
-
=
\int 
\int \int 
\int 
\int \int 
(
)
(
)
z
z
dz
d
z
d
dz
d
z
d
\pi 
\pi 
\pi 
\theta 
\rho \rho \rho 
\theta 
\rho 
\rho \rho 
-
+
-
=
\int 
\int 
\int 
\int 
\int 
\int 
. 
\textcolor{lyred}{注}.
(
)
(
)
x
y
x
y
x
y
x
y
I
dxdy
x
y
z dz
dxdy
z
x
y
dz
+
+
\le 
+
\le 
+
=
+
-
+
-
+
=
\int \int 
\int 
\int \int 
\int 
(
)
(
)
d
d
z dz
d
d
z
dz
\rho 
\pi 
\pi 
\rho 
\pi 
\theta \rho \rho 
\rho 
\theta \rho \rho 
\rho 
-
+
-
=
\int 
\int 
\int 
\int 
\int 
\int 
. 
\textcolor{lyred}{例}.求
(
)
I
x
y
z
dv
\Omega 
=
+
+
\int \int \int 
,其中\Omega 由
z
x
y
=
+
,
z
x
y
=
-
-
围成. 
解.
(
)(
)
z
x
y
z
z
z
z
z
z
x
y
=
+

\Rightarrow 
=
-
\Rightarrow 
+
-
=
\Rightarrow 
=
=
-
-

,即
x
y
+
=, 
故
(
)
(
)
I
x
y
z
dv
x
y
dv
z dv
I
I
\Omega 
\Omega 
\Omega 
=
+
+
=
+
+
=
+
\int \int \int 
\int \int \int 
\int \int \int 
,而 
(
)
x
y
x
y
x
y
I
dxdy
x
y
dz
d
d
dz
\rho 
\pi 
\rho 
\theta \rho \rho 
\rho 
\pi 
\pi 
-
-
-
+
\le 
+
=
+
=
=
-
\int \int 
\int 
\int 
\int 
\int 
, 
(
)
x
y
z
x
y
z
I
dz
z dxdy
dz
z dxdy
z
zdz
z
z
dz
\pi 
\pi 
+
\le 
+
\le -
=
+
=
\cdots 
+
\cdots 
-
=
\int 
\int \int 
\int 
\int \int 
\int 
\int 
\pi 
\pi 
-
,故
I
\pi 
\pi 
=
-
. 
{\bfseries \textcolor{lyred}{3.交换积分顺序}}

\textcolor{lyred}{例}.证明:
()
(
)
()
y
x
dx dy
f z dz
z
f z dz
=
-
\int \int \int 
\int 
. 
证.
()
()
(
)
()
y
x
x
x
x
z
dx dy
f z dz
dx dz
f z dy
dx
x
z f z dz
=
=
-
=
\int \int \int 
\int \int \int 
\int \int 
(
)
()
()
(
)
()
z
z
x
dz
x
z f z dx
zx
f z dz
z
f z dz


-
=
-
=
-




\int \int 
\int 
\int 
,证毕. 
\textcolor{lyred}{例}.求
x
y
I
dx dy y
z dz
=
+
\int \int \int 
. 
解.
z
z
z
z
z
x
x
x
x
I
dx dz y
z dy
dz dx y
z dy
z dz dx ydy
=
+
=
+
=
+
=
\int \int \int 
\int \int \int 
\int 
\int \int 
(
)
2 2
z z
x
z
z
z dz
dx
z dz
z
z dz


-
+
=
-
+
=
+
=
-




\int 
\int 
\int 
\int 
. 
\vspace{4pt}
{\bfseries \textcolor{lyred}{五.柱面坐标下的计算}}

空间点
(
)
, ,
M x y z ,若它在xOy 面上投影的极坐标为(
)
,
\theta \rho ,则称(
)
,
, z
\theta \rho 
为M 的 
\textcolor{lyred}{柱面坐标},并且有关系
cos
sin
x
y
z
z
\rho 
\theta 
\rho 
\theta 
=


=

=

. 
若
()
()
(
)
(
)
:
,
,
,
,
z
z
z
\alpha 
\theta 
\beta \rho \theta 
\rho 
\rho 
\theta 
\theta \rho 
\theta \rho 
\Omega 
\le 
\le 
\le 
\le 
\le 
\le 
,则 
(
)
()
()
(
)
(
)
(
)
,
,
, ,
cos ,
sin ,
z
z
f x y z dv
d
d
f
z dz
\rho 
\theta 
\theta \rho 
\beta 
\alpha 
\rho \theta 
\theta \rho 
\theta 
\rho \rho 
\rho 
\theta \rho 
\theta 
\Omega 
=
\int \int \int 
\int 
\int 
\int 
. 
\textcolor{lyred}{例}.(1)设
:
x
y
z
\Omega 
+
+
\le ,则
(
)
, ,
I
f x y z dv
\Omega 
=
=
\int \int \int 
(
)
(
)
, ,
cos ,
sin ,
x
y
x
y
x
y
dxdy
f x y z dz
d
d
f
z dz
\rho 
\pi 
\rho 
\theta \rho \rho 
\rho 
\theta \rho 
\theta 
-
-
-
+
\le 
-
-
-
-
-
=
\int \int 
\int 
\int 
\int 
\int 
. 
(2)设\Omega 由
x
y
x
+
=
,
z =
,
z =围成,则
(
)
, ,
I
f x y z dv
\Omega 
=
=
\int \int \int 
(
)
(
)
2cos
, ,
cos ,
sin ,
x
y
x
dxdy
f x y z dz
d
d
f
z dz
\pi 
\theta 
\pi 
\theta 
\rho \rho 
\rho 
\theta \rho 
\theta 
+
\le 
-
=
\int \int 
\int 
\int 
\int 
\int 
. 
\textcolor{lyred}{例}.设(
)
(
)
,0
,
,
x
y
r
z h
u r h
f
x
y
z dv
+
\le 
\le \le 
=
+
\int \int \int 
,其中f 连续,求
2u
r h
\partial 
\partial \partial . 
解. (
)
(
)
(
)
,
,
,
r
h
r
h
u r h
d
d
f
z dz
d
f
z dz
\pi 
\theta \rho \rho 
\rho 
\pi \rho \rho 
\rho 
=
=
\int 
\int 
\int 
\int 
\int 
,故 
(
)
,
h
u
r
f r z dz
r
\pi 
\partial 
=
\partial 
\int 
,
(
)
,
u
rf r h
r h
\pi 
\partial 
=
\partial \partial 
. 
\textcolor{lyred}{注}. (
)
(
)
(
)
,
,
,
h
r
h
r
u r h
dz
d
f
z
d
dz
f
z
d
\pi 
\theta 
\rho 
\rho \rho 
\pi 
\rho 
\rho \rho 
=
=
\int 
\int 
\int 
\int \int 
,故 
(
)
,
r
u
f
h
d
h
\pi 
\rho 
\rho \rho 
\partial 
=
\partial 
\int 
,
(
)
,
u
rf r h
h r
\pi 
\partial 
=
\partial \partial 
. 
\vspace{4pt}
{\bfseries \textcolor{lyred}{六.球面坐标下的计算}}

空间点M 的经度,纬度以及到原点的距离构成有序数组(
)
, ,r
\theta \varphi 
,它称为该点的 
\textcolor{lyred}{球面坐标},有关系
sin
cos
sin
sin
cos
x
r
y
r
z
r
\varphi 
\theta 
\varphi 
\theta 
\varphi 
=


=

=

. 
若
(
)
(
)
:
,
,
,
,
r
r
r
\theta 
\theta 
\theta \varphi 
\varphi 
\varphi 
\varphi \theta 
\varphi \theta 
\Omega 
\le 
\le 
\le 
\le 
\le 
\le 
,则 
(
)
(
)
(
)
(
)
,
,
, ,
sin
sin
cos , sin
sin , cos
r
r
f x y z dv
d
d
f r
r
r
r dr
\varphi \theta 
\theta 
\varphi 
\theta 
\varphi 
\varphi \theta 
\theta 
\varphi \varphi 
\varphi 
\theta 
\varphi 
\theta 
\varphi 
\Omega 
=
\int \int \int 
\int 
\int 
\int 
. 
\textcolor{lyred}{例}.求
(
)
I
x
y
z
dv
\Omega 
=
+
+
\int \int \int 
,其中
:
x
y
z
\Omega 
+
+
\le . 
解.
(
)
(
)
I
x
y
z
dv
x dv
x
y
z
dv
\Omega 
\Omega 
\Omega 
=
+
+
=
=
+
+
=
\int \int \int 
\int \int \int 
\int \int \int 
sin
sin
d
d
r
r dr
d
r dr
\pi 
\pi 
\pi 
\theta 
\varphi \varphi 
\pi 
\varphi \varphi 
\pi 
\cdots 
=
\cdots 
\cdots 
\cdots 
=
\int 
\int 
\int 
\int 
\int 
. 
\textcolor{lyred}{例}.求
I
x
y
z
dv
\Omega 
=
+
+
-
\int \int \int 
,其中
:
x
y
z
x
y
\Omega 
+
\le 
\le 
-
-
. 
解.
(
)
(
)
I
x
y
z
dv
x
y
z
dv
\Omega 
\Omega 
=
-
+
+
+
+
+
-
=
\int \int \int 
\int \int \int 
(
)
(
)
(
)
sin
sin
d
d
r r dr
d
d
r
r dr
\pi 
\pi 
\pi 
\pi 
\pi 
\theta 
\varphi \varphi 
\theta 
\varphi \varphi 
-
-
+
-
=
\int 
\int 
\int 
\int 
\int 
\int 
. 
\textcolor{lyred}{例}.求
(
)
I
x
y
z
dv
\Omega 
=
+
+
\int \int \int 
,其中
: x
y
z
x
y
z
\Omega 
+
+
\le 
+
+
. 
解.
:
x
y
z






\Omega 
-
+
-
+
-
\le 












,令
,
,
x
x
y
y
z
z
'
'
'
=
+
=
+
=
+
, 
则
I
x
y
z
dv
x
y
z
dv
'
'
\Omega 
\Omega 










'
'
'
'
'
'
'
'
=
+
+
+
+
+
=
+
+
+
=






















\int \int \int 
\int \int \int 
3 4
3 3
sin
4 3
d
d
r
r dr
\pi 
\pi 
\theta 
\varphi \varphi 
\pi 
\pi 


\cdots 
+
\cdots 
=






\int 
\int 
\int 
. 
\textcolor{lyred}{例}.求
(
)
I
x
y
dv
\Omega 
=
+
\int \int \int 
,其中
(
)
: x
y
z
a
a
\Omega 
+
+
-
\le 
. 
解.
(
)
(
)
2 cos
sin
sin
a
I
x
y
dv
x
y
dv
d
d
r
r dr
\pi 
\varphi 
\pi 
\theta 
\varphi \varphi 
\varphi 
\Omega 
\Omega 
=
+
=
+
=
\cdots 
=
\int \int \int 
\int \int \int 
\int 
\int 
\int 
sin
cos
a
d
a
\pi 
\pi 
\varphi 
\varphi \varphi 
\pi 
=
\int 
. 
\textcolor{lyred}{例}.求半径a 的球面与半顶角为\alpha 的内接圆锥面所围立体体积. 
解.
(
)
2 cos
sin
1 cos
a
a
V
dv
d
d
r dr
\varphi 
\pi 
\alpha 
\pi 
\theta 
\varphi \varphi 
\alpha 
\Omega 
=
=
=
-
\int \int \int 
\int 
\int 
\int 
. 
\textcolor{lyred}{例}.求
I
z dv
\Omega 
=\int \int \int 
,其中
:
z
x
y
z
z
\Omega 
\le 
+
+
\le 
. 
解.
(
)
2cos
cos
sin
cos
I
d
d
r
r dr
\pi 
\varphi 
\pi 
\varphi 
\theta 
\varphi \varphi 
\varphi 
\pi 
=
\cdots 
=
\int 
\int 
\int 
. 
\textcolor{lyred}{例}.求
I
z dv
\Omega 
=\int \int \int 
,其中
:
x
y
z
\Omega 
+
+
\le ,
x
y
z
z
+
+
\le 
. 
解.
x
y
z
z
x
y
z
z

+
+
=
\Rightarrow 
=

+
+
=

,
x
y
+
=
,在两曲面交线处
\pi 
\varphi =
,故 
(
)
(
)
2cos
sin
cos
sin
cos
I
d
d
r
r dr
d
d
r
r dr
\pi 
\pi 
\varphi 
\pi 
\pi 
\pi 
\theta 
\varphi \varphi 
\varphi 
\theta 
\varphi \varphi 
\varphi 
\pi 
=
+
=
\int 
\int 
\int 
\int 
\int 
\int 
. 
\textcolor{lyred}{补充练习 }
1.将
(
)
, ,
x y
x
dx
dy
f x y z dz
+
-
\int 
\int 
\int 
化为先x ,再z ,后y 的三次积分. 
解.
(
)
(
)
, ,
, ,
x y
y
x y
x
dx
dy
f x y z dz
dy
dx
f x y z dz
+
-
+
-
=
=
\int 
\int 
\int 
\int 
\int 
\int 
(
)
(
)
, ,
, ,
y
y
y
y
z y
dy dz
f x y z dx
dy dz
f x y z dx
-
-
-
+
\int \int 
\int 
\int \int 
\int 
. 
2.两种方法计算
I
z dv
\Omega 
=\int \int \int 
,其中
:
x
y
z
\Omega 
+
+
\le ,
z
x
y
+\ge 
+
. 
解.
x
y
x
y
x
y
I
dxdy
z dz
d
d
z dz
\rho 
\pi 
\rho 
\pi 
\theta \rho \rho 
-
-
-
-
+
\le 
+
-
=
=
=
\int \int 
\int 
\int 
\int 
\int 
; 
(
)
x
y
z
x
y
z
I
dz
z dxdy
dz
z dxdy
\pi 
-
+
\le -
+
\le 
+
=
+
=
\int 
\int \int 
\int 
\int \int 
. 
3.两种方法计算
z
x
y
=
+
与
z
x
y
=
-
-
所围立体的体积. 
解.
(
)
x
y
x
y
x
y
V
dv
dxdy
dz
d
d
\pi 
\theta 
\rho 
\rho \rho \rho 
\pi 
-
-
\Omega 
+
\le 
+
=
\cdots 
=
\cdots 
=
-
-
=
\int \int \int 
\int \int 
\int 
\int 
\int 
; 
(
)
x
y
z
x
y
z
V
dv
dz
d
dz
d
z dz
z dz
\sigma 
\sigma 
\pi 
\pi 
\pi 
\Omega 
+
\le 
+
\le -
=
\cdots 
=
+
=
+
-
=
\int \int \int 
\int 
\int \int 
\int 
\int \int 
\int 
\int 
. 
4.两种方法计算I
zdv
\Omega 
=\int \int \int 
,其中
:
x
y
z
\Omega 
+
+
\le ,
x
y
z
z
+
+
\le 
. 
解.
x
y
z
z
z
x
y
z
z

+
+
=
\Rightarrow =
\Rightarrow 
=

+
+
=

,即
x
y
+
=
,故 
(
)
3 4
3 4
2 1
x
y
x
y
x
y
x
y
I
dxdy
zdz
x
y
dxdy
\pi 
-
-
+
\le 
+
\le 
-
-
-
=
=
-
-
-
=
\int \int 
\int 
\int \int 
; 
1 2
1 2
x
y
z z
x
y
z
I
dz
zdxdy
dz
zdxdy
\pi 
+
\le 
-
+
\le -
=
+
=
\int 
\int \int 
\int 
\int \int 
. 
5.设
()
f x 连续,
()
lim
x
f x
x
\to 
=,求
(
)
lim
t
x
y
z
t z
I
f x
y
z
dv
t
+
\to 
+
+
\le 
\cdots 
=
+
+
\int \int \int 
. 
解.
(
)
()
lim
t
I
f
t
f
t
\psi 
\eta 

\pi 
\pi 
+
\to 
=
+
+
\cdots 
=
=
. 
6.设
:
t
D
x
y
t
+
\le 
,
:
t x
y
z
t
\Omega 
+
+
\le 
,
()
(
)
0,
f x
C
\in 
+\infty ,且
()
f x >
,求 
(
)
(
)
lim
t
t
t
D
f x
y
z
dv
t
f x
y
d\sigma 
+
\Omega 
\to 
+
+
+
\int \int \int 
\int \int 
. 
解.原式
(
)
(
)
()
()
lim
t
f
t
f
f
f u
v
t
\psi 

\eta 
\pi 
\pi 
+
\to 
+
+
\cdots 
=
=
\cdots 
=
+
\cdots 
. 
\vspace{6pt}
{\large \bfseries \textcolor{lyred}{第10.4 节 重积分的应用}}

\vspace{4pt}
{\bfseries \textcolor{lyred}{一.曲面的面积}}

\textcolor{lyred}{定理}.设光滑曲面
(
)
:
,
z
z x y
\Sigma 
=
,(
)
,
xy
x y
D
\in 
,其中
xy
D 为有界闭区域,则它的面积 
xy
x
y
D
A
z
z d\sigma 
=
+
+
\int \int 
. 
\textcolor{lyred}{例}.求半径为R 的球的表面积. 
解.
z
R
x
y
=
-
-
,故
x
y
R
R
A
d
R
R
x
y
\sigma 
\pi 
+
\le 
=
=
-
-
\int \int 
. 
\textcolor{lyred}{例}.求
x
y
z
a
+
+
=
含在
x
y
ax
+
=
内部分的面积. 
解.
(
)
cos
a
x
y
ax
a
a
A
d
d
d
a
a
x
y
a
\pi 
\theta 
\pi 
\sigma 
\theta 
\rho \rho 
\pi 
\rho 
+
\le 
-
=
=
\cdots 
=
-
-
-
-
\int \int 
\int 
\int 
. 
\textcolor{lyred}{例}.求
z
x
y
=
+
被
z
x
=
所割下部分的面积. 
解.
z
x
y
x
y
x
z
x
=
+

\Rightarrow 
+
=

=

,故
2 2
x
y
x
A
d\sigma 
\pi 
+
\le 
=
\cdots 
=
\int \int 
. 
\textcolor{lyred}{例}.求
1 : z
x
y
\Sigma 
=
+
与
2 :
z
x
y
\Sigma 
=
-
+
围成立体的表面积. 
解.
z
x
y
x
y
z
x
y
=
+

\Rightarrow 
+
=
=
-
+

,故
A
A
A
=
+
= 
5 5
x
y
x
y
d
x
y d
\sigma 
\sigma 
\pi 
+
\le 
+
\le 


-
+
+
+
=
+






\int \int 
\int \int 
. 
\textcolor{lyred}{例}.某屋顶由两个半径分别为1和2 的半球叠加而成,求表面积. 
解.
4 3
x
y
A
d
x
y
\pi 
\sigma 
\pi 
\pi 
\le 
+
\le 
=
+
=
+
-
-
\int \int 
. 
\vspace{4pt}
{\bfseries \textcolor{lyred}{二.质心与形心}}

{\bfseries \textcolor{lyred}{1.质点系}}

设xOy 平面上有n 个位于(
)
(
)
,
,
,
,
n
n
x y
x
y
\Lambda 
处的质点,质量为
1,
,
n
m
m
\Lambda 
,则 
该质点系的\textcolor{lyred}{质心}坐标为(
)
,x y ,其中 
n
i
i
n
y
i
i
i
n
i
i
i
m x
M
m
x
x
M
M
m
=
=
=
=
=
=
\sum 
\sum 
\sum 
,
n
i
i
n
x
i
i
i
n
i
i
i
m y
M
m
y
y
M
M
m
=
=
=
=
=
=
\sum 
\sum 
\sum 
. 
若
i
m
m
=
,则
n
i
i
x
x
n
=
=\sum 
,
n
i
i
y
y
n
=
=\sum 
,此时(
)
,x y 也称为质点系的\textcolor{lyred}{形心}. 
{\bfseries \textcolor{lyred}{2.平面薄片}}

设平面薄片占有xOy 面上有界闭区域D ,密度为
(
)
,x y
\mu 
,连续,则它的\textcolor{lyred}{质心}坐标\textcolor{lyred}{ }
(
)
(
)
,
,
y
D
D
x
x y d
M
x
M
x y d
\mu 
\sigma 
\mu 
\sigma 
=
=\int \int 
\int \int 
,
(
)
(
)
,
,
x
D
D
y
x y d
M
y
M
x y d
\mu 
\sigma 
\mu 
\sigma 
=
=\int \int 
\int \int 
. 
若薄片均匀,则
D
D
xd
x
d
\sigma 
\sigma 
=\int \int 
\int \int 
,
D
D
yd
y
d
\sigma 
\sigma 
=\int \int 
\int \int 
,此时(
)
,x y 也称为D 的\textcolor{lyred}{形心}. 
\textcolor{lyred}{例}.求位于两圆
2sin
\rho 
\theta 
=
和
4sin
\rho 
\theta 
=
之间的均匀薄片的质心. 
解.
x =
,
4sin
2sin
sin
D
y
yd
d
d
A
\pi 
\theta 
\theta 
\sigma 
\theta 
\rho 
\theta \rho \rho 
\pi 
=
=
\cdots 
=
\int \int 
\int 
\int 
,故为
0, 3






. 
\textcolor{lyred}{例}.在均匀半圆形薄板下接一个宽与圆直径相同的长方形的均匀薄板,要求质心 
在圆心处,求圆半径与长方形的高之比. 
解.以半圆直径边为x 轴,对称轴为y 轴,则质心为原点,故
y =
, 
于是
sin
R
R
x
D
D
R
h
M
yd
yd
d
d
dx
ydy
\pi 
\mu 
\sigma 
\mu 
\sigma 
\mu 
\theta \rho 
\theta \rho \rho 
\mu 
-
-
=
+
=
\cdots 
+
=
\Rightarrow 
\int \int 
\int \int 
\int 
\int 
\int 
\int 
:
:
R
Rh
R h
\mu 
\mu 
\mu 
\mu 
-
=
\Rightarrow 
=
. 
{\bfseries \textcolor{lyred}{3.空间立体}}

设物体占有空间有界闭区域\Omega ,密度为
(
)
, ,
x y z
\rho 
,连续,则\textcolor{lyred}{质心}坐标 
x dv
x
dv
\rho 
\rho 
\Omega 
\Omega 
=\int \int \int 
\int \int \int 
,
y dv
y
dv
\rho 
\rho 
\Omega 
\Omega 
=\int \int \int 
\int \int \int 
,
z dv
z
dv
\rho 
\rho 
\Omega 
\Omega 
=\int \int \int 
\int \int \int 
. 
若\Omega 均匀,则
xdv
x
dv
\Omega 
\Omega 
=\int \int \int 
\int \int \int 
,
ydv
y
dv
\Omega 
\Omega 
=\int \int \int 
\int \int \int 
,
zdv
z
dv
\Omega 
\Omega 
=\int \int \int 
\int \int \int 
,此时(
)
, ,
x y z 也称为\Omega 的\textcolor{lyred}{形心}. 
\textcolor{lyred}{例}.求均匀半球体
(
)
:
x
y
z
a
z
\Omega 
+
+
\le 
\ge 
的质心. 
解.
x
y
=
=
,
a
x
y
a
z
a
z
zdv
dz
zdxdy
a
V
V
a
\pi 
\pi 
\Omega 
+
\le 
-
=
=
=
\cdots 
=
\int \int \int 
\int 
\int \int 
. 
\textcolor{lyred}{例}.均匀半球,密度为
\rho ,下接半径相同的均匀圆柱体,密度为
\rho ,质心在球心处, 
求圆柱半径与高之比. 
解.以球心为原点,对称轴为z 轴建立坐标系,设质心(
)
, ,
x y z ,由 
z
z dv
zdv
zdv
\rho 
\rho 
\rho 
\Omega 
\Omega 
\Omega 
=
\Rightarrow 
=
\Rightarrow 
+
=
\int \int \int 
\int \int \int 
\int \int \int 
,于是 
sin
cos
:
R
h
x
y
R
d
d
r
r dr
dxdy
zdz
R h
\pi 
\pi 
\rho 
\rho 
\theta 
\varphi \varphi 
\varphi 
\rho 
\rho 
-
+
\le 
\cdots 
+
=
\Rightarrow 
=
\int 
\int 
\int 
\int \int 
\int 
. 
\vspace{4pt}
{\bfseries \textcolor{lyred}{三.转动惯量(惯性矩)}}

{\bfseries \textcolor{lyred}{1.质点系}}

设xOy 平面上有n 个质点,分别位于(
)
(
)
,
,
,
,
n
n
x y
x
y
\Lambda 
处,则该质点系对于x 轴及 
y 轴的\textcolor{lyred}{转动惯量}分别为
n
x
i
i
i
I
m y
=
=\sum 
,
n
y
i
i
i
I
m x
=
=\sum 
. 
{\bfseries \textcolor{lyred}{2.平面薄片}}

设平面薄片占有xOy 面上的有界闭区域D ,密度为
(
)
,x y
\mu 
,连续,则它对于x 轴
和y 轴的\textcolor{lyred}{转动惯量}
(
)
,
x
D
I
y
x y d
\mu 
\sigma 
=\int \int 
,
(
)
,
y
D
I
x
x y d
\mu 
\sigma 
=\int \int 
. 
\textcolor{lyred}{例}.求半径为a ,密度为\mu 的半圆薄片对其直径边的转动惯量. 
解.
,
sin
a
x
x
y
a
y
I
y d
d
d
a
\pi 
\mu 
\sigma 
\mu 
\theta \rho 
\theta \rho \rho 
\pi \mu 
+
\le 
\ge 
=
=
\cdots 
=
\int \int 
\int 
\int 
. 
\textcolor{lyred}{例}.求
y
x
=
,
y =围成的密度为\mu 的薄片对
y =-的转动惯量. 
解.
(
)
(
)
D
x
I
y
d
dx
y
dy
\mu \sigma 
\mu 
\mu 
-
=
+
=
+
=
\int \int 
\int 
\int 
. 
{\bfseries \textcolor{lyred}{3.空间立体}}

设物体占有空间有界闭区域\Omega ,密度
(
)
, ,
x y z
\rho 
,连续,则转动惯量为 
(
)
xI
y
z
dv
\rho 
\Omega 
=
+
\cdots 
\int \int \int 
,
(
)
yI
z
x
dv
\rho 
\Omega 
=
+
\cdots 
\int \int \int 
,
(
)
zI
x
y
dv
\rho 
\Omega 
=
+
\cdots 
\int \int \int 
. 
\textcolor{lyred}{例}.设匀质立体
(
)
:
1 1
x
y
z
z
\Omega 
+
-
\le 
\le 
\le 
的密度为
\rho ,求
zI . 
解.
(
)
(
)
z
x
y
z
I
x
y
dv
dz
x
y
dxdy
\rho 
\rho 
\pi \rho 
\Omega 
+
\le 
+
=
+
=
+
=
\int \int \int 
\int 
\int \int 
. 
\textcolor{lyred}{例}.求密度为
\rho 的均匀球体关于过球心的轴的转动惯量. 
解.
(
)
(
)
zI
x
y
dv
x
y
z
dv
\rho 
\rho 
\Omega 
\Omega 
=
+
=
+
+
=
\int \int \int 
\int \int \int 
sin
a
d
d
r
r dr
a
\pi 
\pi 
\rho 
\theta 
\varphi \varphi 
\pi \rho 
\cdots 
=
\int 
\int 
\int 
. 
\vspace{4pt}
{\bfseries \textcolor{lyred}{四.引力}}

设物体占有空间有界闭域\Omega ,密度为
(
)
, ,
x y z
\rho 
.连续,则它对
(
)
,
,
P x
y z
处单位 
质点的\textcolor{lyred}{引力}为 
(
)
(
)
(
)
(
)
,
,
,
,
x
y
z
x
x
y
y
z
z
F
F F F
G
dv
dv
dv
r
r
r
\rho 
\rho 
\rho 
\Omega 
\Omega 
\Omega 


-
-
-
=
=




\int \int \int 
\int \int \int 
\int \int \int 
\rho 
,其中 
(
)
(
)
(
)
r
x
x
y
y
z
z
=
-
+
-
+
-
. 
\textcolor{lyred}{例}.求密度为
\rho 的球体
: x
y
z
R
\Omega 
+
+
\le 
对位于点
(
)(
)
0 0,0,
M
a
a
R
>
处的单位 
质点的引力. 
解.由对称性,设
(
)
0,0,
z
F
F
=
\rho 
,则
(
)
(
)
z
z
a
F
G
dv
x
y
z
a
\rho 
\Omega 
-
=
=


+
+
-


\int \int \int 
(
)
(
)
R
R
x
y
R
z
G
z
a dz
dxdy
x
y
z
a
\rho 
-
+
\le 
-
-
=


+
+
-


\int 
\int \int 
(
)
(
)
R
R
z
R
R
d
GM
G
z
a dz
d
G
a
a
z
a
\pi 
\pi 
\rho 
\rho \rho 
\rho 
\theta 
\rho 
-
-
-
=-
=-


+
-


\int 
\int 
\int 
. 
\textcolor{lyred}{例}.设密度为
\rho 的均匀柱体
(
)
:
x
y
R
z
h
\Omega 
+
\le 
\le \le 
,求它对位于
(
)(
)
0,0,
M
a a
h
>
处的单位质点的引力. 
解.设
(
)
0,0,
z
F
F
=
\rho 
,则
(
)
(
)
(
)
(
)
z
z
a
F
G
dv
x
y
z
a
\rho 
\Omega 
-
=
=


-
+
-
+
-


\int \int \int 
(
)
(
)
(
)
(
)
h
R
h
x
y
R
z
a
z
a
G
dxdy
dz
G
d
d
dz
x
y
z
a
z
a
\pi 
\rho 
\rho 
\theta \rho \rho 
\rho 
+
\le 
-
-
=
=




+
+
-
+
-




\int \int 
\int 
\int 
\int 
\int 
(
)
2 G
h
R
h
a
R
a
\pi 
\rho 

-
+
+
-
-
+



. 
\textcolor{lyred}{例}.求半径为R ,半顶角为\alpha ,密度为
\rho 的均匀球锥体对顶点处单位质点的引力. 
解.取顶点为原点,对称轴为z 轴建立坐标系,则
(
)
0,0,
z
F
F
=
\rho 
,其中 
(
)
cos
sin
sin
R
z
zdv
r
F
G
G
d
d
r
dr
G
R
r
x
y
z
\pi 
\alpha 
\rho 
\varphi 
\rho 
\theta 
\varphi 
\varphi 
\rho \pi 
\alpha 
\Omega 
=
=
\cdots 
=
+
+
\int \int \int 
\int 
\int 
\int 
. 
\textcolor{lyred}{例}.设顶点在原点,半顶角为6
\pi 的球面
(
)
x
y
z
a
a
+
+
-
=
的内接均匀正圆锥体 
\Omega (密度为
\rho )对原点处单位质点的引力. 
解.
x
y
F
F
=
=
,
z
Gz
dv
F
r
\rho 
\Omega 
=\int \int \int 
,\Omega 顶部高度
2 cos
cos
z
a
a
\pi 
\pi 


=
=




,于是 
在球坐标中,\Omega 的顶面方程为
cos
r
a
\varphi =
,即
2cos
a
r
\varphi 
=
,因此 
2cos
cos
sin
a
z
r
F
G
d
d
r dr
aG
r
\pi 
\varphi 
\pi 
\varphi 
\rho 
\theta 
\varphi \varphi 
\rho 
\pi 


=
\cdots 
=
-






\int 
\int 
\int 
. 
\textcolor{lyred}{例}.设有密度为\mu 的均匀半圆形薄片,占有区域
:
D x
y
R
+
\le 
,
y \ge 
,圆心上方 
有一单位质点P , OP
a
=
,求薄片对该质点的引力. 
解.质点位于(
)
0,0,a ,设
(
)
0,
,
y
z
F
F F
=
\rho 
,其中
(
)
y
D
y
F
G
d
r
\mu 
\sigma 
-
=
=
\int \int 
(
)
sin
ln
R
d
R
a
R
R
G
d
G
a
a
R
a
\pi 
\rho 
\theta \rho \rho 
\mu 
\theta 
\mu 
\rho 


\cdots 
+
+
=
-




+


+
\int 
\int 
, 
(
)
(
)
R
z
D
a
d
a
F
G
d
G a d
G
r
a
R
a
\pi 
\mu 
\rho \rho 
\sigma 
\mu 
\theta 
\mu \pi 
\rho 


-
=
=-
=-
-


+


+
\int \int 
\int 
\int 
. 